% !TEX root = ../main.tex
\chapter{Zoo Architecture}
\label{ch:zoo_architecture}

\Lucid\ ships a model zoo: a curated collection of architectures
under \code{lucid.models}, organised into three domains (vision,
text, generative), with consistent API conventions, pre-trained
weight loading via a hub, and a structured family contract that
every model follows. The zoo serves three purposes:

\begin{itemize}
  \item \textbf{Reproducibility}: the standard architectures
    (ResNet, ViT, BERT, GPT-2, DDPM) are implemented exactly as
    in the original papers, providing a baseline that user code
    can match.
  \item \textbf{Benchmarking}: the zoo entries are the
    canonical workloads for measuring \Lucid\ performance on
    Apple silicon.
  \item \textbf{Building blocks}: users compose new architectures
    from zoo components (a custom transformer reusing the BERT
    block, a custom detection head on a ResNet backbone).
\end{itemize}

This chapter describes the zoo's organisational structure: the
three-domain split, the per-family file layout, the
\code{@register\_model} decorator, the \code{@model\_family\_meta}
metadata, the \code{AutoModel} factory, and the contract that
Hard Rules H10--H12 enforce.

\section{The Three Domains}
\label{sec:zoo_architecture:three_domains}

\code{lucid.models} is partitioned into three subpackages:

\begin{itemize}
  \item \code{lucid.models.vision}: convolutional networks (ResNet,
    VGG, DenseNet, MobileNet, EfficientNet, ConvNeXt), vision
    transformers (ViT, Swin, MaxViT, PVT), object detection
    (Faster R-CNN, Mask R-CNN, DETR, YOLO), segmentation
    (U-Net). Roughly 62 model families.
  \item \code{lucid.models.text}: language models (Transformer
    encoder-decoder, BERT, GPT, GPT-2, RoFormer). Roughly 7
    families.
  \item \code{lucid.models.generative}: generative models (VAE,
    DDPM, NCSN). 3 families.
\end{itemize}

The grouping is by primary use case rather than by
architectural family (a CNN-based detector lives under vision,
not under CNN, even though it uses convolution).

\section{Per-Family File Layout}
\label{sec:zoo_architecture:file_layout}

Every model family follows the same five-slot file structure:

\begin{lstlisting}
lucid/models/<domain>/<family>/
    __init__.py        # public API exports
    _config.py         # Config dataclass
    _model.py          # nn.Module subclasses
    _pretrained.py     # @register_model factories
\end{lstlisting}

\subsection{The Config dataclass}
\label{ssec:zoo_architecture:config_dataclass}

Each family's \code{\_config.py} defines a frozen dataclass that
parameterises the family:

\begin{lstlisting}[style=lucid,language=LucidPython]
from dataclasses import dataclass
from typing import ClassVar
from lucid.models import model_family_meta

@model_family_meta(
    canonical_name="ResNet",
    citation="He et al., 2015 (arXiv:1512.03385)",
    theory=r"""
        Deep residual learning. The key building block is the
        residual unit y = F(x) + x, which lets gradient flow
        directly through the skip connection in addition to
        through F. This addresses the degradation problem in
        very deep networks.
    """,
)
@dataclass(frozen=True)
class ResNetConfig:
    model_type: ClassVar[str] = "resnet"
    layers: tuple[int, ...] = (3, 4, 6, 3)  # ResNet-50 default
    num_classes: int = 1000
    in_channels: int = 3
    dropout: float = 0.0
\end{lstlisting}

\code{frozen=True} makes the config immutable: configs are passed
by value, never mutated after construction. The
\code{model\_type} class variable is the registry key.

\subsection{The Model class}
\label{ssec:zoo_architecture:model_class}

\code{\_model.py} defines the actual \code{nn.Module} subclass:

\begin{lstlisting}[style=lucid,language=LucidPython]
import lucid.nn as nn
from ._config import ResNetConfig

class ResNet(nn.Module):
    def __init__(self, config: ResNetConfig):
        super().__init__()
        self.config = config
        self.stem = nn.Conv2d(config.in_channels, 64, ...)
        self.layer1 = self._make_layer(...)
        # ...

    def forward(self, x):
        x = self.stem(x)
        # ...
        return x
\end{lstlisting}

The Direct model class takes the Config in its constructor. Task
wrappers (\code{ResNetForImageClassification}) wrap the Direct
model with the appropriate output head.

\subsection{The Pretrained factories}
\label{ssec:zoo_architecture:pretrained_factories}

\code{\_pretrained.py} contains the \code{@register\_model}
decorated factory functions:

\begin{lstlisting}[style=lucid,language=LucidPython]
from lucid.models import register_model
from ._config import ResNetConfig
from ._model import ResNet

@register_model(
    task="image_classification",
    family="resnet",
    model_type="resnet",
    model_class=ResNet,
    default_config=ResNetConfig(layers=(3, 4, 6, 3)),
    params=25_557_032,    # paper-cited
    summary="auto",
)
def resnet_50(**kwargs):
    config = ResNetConfig(layers=(3, 4, 6, 3), **kwargs)
    return ResNet(config)
\end{lstlisting}

The decorator registers the factory in a global model registry,
which the AutoModel API and the docs site introspect.

\section{The H10 Variant Rule}
\label{sec:zoo_architecture:h10_variant_rule}

Hard Rule H10 prescribes the naming convention for size
variants: full words, not abbreviations.

\begin{itemize}
  \item Correct: \code{resnet\_50}, \code{vit\_base\_16},
    \code{swin\_tiny}.
  \item Incorrect: \code{resnet50}, \code{vit\_b\_16},
    \code{swin\_t}.
\end{itemize}

The full-word convention is verbose but unambiguous and
search-friendly. The Hard Rule applies to factory names,
registry entries, test ids, and any documentation reference.

\section{The H11 Paper-Cited Variants Rule}
\label{sec:zoo_architecture:h11_paper_rule}

Hard Rule H11 prescribes that only variants explicitly listed in
the originating paper can be registered as factories. For
example:

\begin{itemize}
  \item \code{vit\_base\_16}: in \cite{dosovitskiy_vit_2021}'s
    Table 1.
  \item \code{vit\_tiny\_16}: not in the paper. Not registered.
\end{itemize}

A user who wants a non-paper-cited size constructs the model
manually via the Direct API:

\begin{lstlisting}[style=lucid,language=LucidPython]
from lucid.models.vision.vit import ViTConfig, ViT

cfg = ViTConfig(hidden_size=96, num_layers=6, num_heads=3, ...)
model = ViT(cfg)
\end{lstlisting}

The rule prevents the registry from accumulating made-up
variants that lack a paper baseline. The exception is for
families where the paper itself lists a small or tiny variant
(YOLOv3-tiny, MobileNetV3-small are paper-blessed).

\section{The H12 Family-Add Procedure}
\label{sec:zoo_architecture:h12_family_add}

Hard Rule H12 prescribes the 16-step procedure for adding a new
model family. The short form:

\begin{enumerate}
  \item Create the four-file structure under
    \code{lucid/models/<domain>/<family>/}.
  \item Define the Config with
    \code{@model\_family\_meta(canonical\_name=, citation=, theory=)}.
  \item Define the Direct model, the task wrappers, the Output
    dataclasses (5-slot model.py structure).
  \item Add factories with
    \code{@register\_model(task=, family=, model\_type=, model\_class=, default\_config=, params=, summary='auto')}.
  \item Run the three-step validation:
    \code{tools.validate\_model\_zoo --family <family> --runtime}
    (0 errors), the family contract pytest, the build-model-
    summaries tool (all factories pass \code{(shadow)}).
\end{enumerate}

Each step has explicit deliverables. A family that does not pass
all three validation tools is rejected at PR review.

\section{AutoModel and AutoConfig}
\label{sec:zoo_architecture:auto}

\code{lucid.models.AutoModel} is the entry point for loading
models by name:

\begin{lstlisting}[style=lucid,language=LucidPython]
from lucid.models import AutoModel

model = AutoModel.from_pretrained("resnet_50")
# Or:
model = AutoModel.from_pretrained("vit_base_16",
                                  num_classes=10)
\end{lstlisting}

The factory:
\begin{enumerate}
  \item Looks up the name in the registry.
  \item Constructs the config (with any kwargs overrides).
  \item Constructs the model.
  \item Downloads + loads pretrained weights from the hub (if
    available).
  \item Returns the configured model.
\end{enumerate}

Task-specific variants: \code{AutoModelForImageClassification},
\code{AutoModelForObjectDetection},
\code{AutoModelForCausalLM}, etc.

\section{The Model Registry}
\label{sec:zoo_architecture:registry}

\code{lucid.models.\_registry} maintains the global registry of
all registered factories. The registry tracks:

\begin{itemize}
  \item Factory function (callable).
  \item Task tag (\code{image\_classification},
    \code{object\_detection}, etc.).
  \item Family name.
  \item Default config.
  \item Paper-cited parameter count.
  \item Auto-generated summary (layer tree).
\end{itemize}

The registry is introspectable:

\begin{lstlisting}[style=lucid,language=LucidPython]
import lucid.models as m

m.list_models()            # ~200 factory names
m.list_models(task="image_classification")
m.is_model("vit_base_16")  # True
m.get_model_info("resnet_50")  # dict of metadata
\end{lstlisting}

\section{The Mixins}
\label{sec:zoo_architecture:mixins}

\Lucid\ provides task-specific mixins under
\code{lucid.models.\_mixins}:

\begin{itemize}
  \item \code{BackboneMixin}: returns feature maps from
    multiple stages (for detection / segmentation heads).
  \item \code{ClassificationHeadMixin}: applies global pooling +
    linear classifier.
  \item \code{GenerationMixin}: provides \code{generate()} for
    causal language models (greedy, beam search, sampling).
  \item \code{DiffusionMixin}: provides forward and reverse
    diffusion utilities for DDPM-family models.
\end{itemize}

A task-wrapper class typically inherits the Direct model class
plus a mixin:

\begin{lstlisting}[style=lucid,language=LucidPython]
class GPT2ForCausalLM(GPT2, GenerationMixin):
    def __init__(self, config):
        super().__init__(config)
        self.lm_head = nn.Linear(config.hidden_size,
                                 config.vocab_size, bias=False)

    def forward(self, input_ids, labels=None):
        hidden = super().forward(input_ids)
        logits = self.lm_head(hidden)
        if labels is not None:
            loss = F.cross_entropy(
                logits.reshape(-1, logits.shape[-1]),
                labels.reshape(-1))
            return CausalLMOutput(loss=loss, logits=logits)
        return CausalLMOutput(logits=logits)
\end{lstlisting}

\section{Output Dataclasses}
\label{sec:zoo_architecture:output_dataclasses}

Each task has a corresponding output dataclass:

\begin{itemize}
  \item \code{ImageClassificationOutput}: logits, loss, hidden
    states, attentions.
  \item \code{ObjectDetectionOutput}: boxes, scores, labels,
    loss components.
  \item \code{CausalLMOutput}: logits, loss, past\_key\_values
    (for KV-cache).
  \item \code{Seq2SeqLMOutput}: encoder/decoder hidden,
    cross-attention, etc.
\end{itemize}

The dataclass returns are structured: the user can access fields
by name (\code{output.logits}, \code{output.loss}) or by tuple
indexing (\code{output[0]}).

\section{Pretrained Weights and the Hub}
\label{sec:zoo_architecture:hub}

The \code{from\_pretrained} factory downloads weights from a hub
(by default, the Hugging Face Hub) and loads them into the
constructed model. The hub URL is in the registry entry; the
download is cached locally.

For weights that are not on a hub, the user passes a local path:

\begin{lstlisting}[style=lucid,language=LucidPython]
model = AutoModel.from_pretrained("/path/to/local/weights.safetensors")
\end{lstlisting}

\chref{ch:pretrained} treats the hub mechanics in detail.

\section{The Validation Tools}
\label{sec:zoo_architecture:validation}

Three tools enforce the family contract.

\subsection{validate\_model\_zoo}
\label{ssec:zoo_architecture:validate}

\code{python -m tools.validate\_model\_zoo --family <family>}
runs static checks:
\begin{itemize}
  \item All factories have the required \code{@register\_model}
    attributes.
  \item The Config has \code{@model\_family\_meta}.
  \item The naming convention (H10) is respected.
  \item Paper-cited parameter counts are present (H11).
\end{itemize}

\code{--runtime} additionally constructs each factory's model
and verifies a one-pass forward succeeds.

\subsection{family contract pytest}
\label{ssec:zoo_architecture:family_pytest}

\code{lucid/test/unit/models/test\_family\_contract.py}
parametrises over all registered families and verifies:
\begin{itemize}
  \item The Config is frozen.
  \item The Direct model and task wrappers exist.
  \item The Output dataclasses exist.
  \item The \code{forward} signature matches the task.
\end{itemize}

\subsection{build\_model\_summaries}
\label{ssec:zoo_architecture:build_summaries}

\code{python -m tools.build\_model\_summaries --family <family>}
runs each factory's shadow allocation (a forward with random
input to introspect the layer tree) and verifies the auto-
generated summary matches the expected structure.

\section{Summary and Forward References}
\label{sec:zoo_architecture:summary}

The model zoo is organised into three domains (vision, text,
generative), with each family following a five-slot file
structure, a Config + Model + Output dataclass + Pretrained
factory layout, and a registry-based discovery mechanism. Hard
Rules H10--H12 prescribe the naming, variant, and family-add
conventions; the three validation tools enforce them at PR
review.

The next chapter, \chref{ch:vision_cnn}, treats the vision CNN
families. \chref{ch:vision_transformers} treats vision
transformers. \chref{ch:detection_segmentation} treats
detection and segmentation. \chref{ch:language_models} treats
language models. \chref{ch:generative} treats generative
models. \chref{ch:pretrained} treats the hub mechanics.

\begin{lucidnote}
For the user: \code{lucid.models.AutoModel.from\_pretrained(name)}
is the standard entry point. The registry name (e.g.,
\code{resnet\_50}) is searchable via \code{list\_models()}.
\end{lucidnote}
